\documentclass[11pt]{article}
\usepackage[letterpaper,hmargin=0.92in,vmargin=0.82in]{geometry}
\usepackage{amsmath,amssymb,amsthm}
\usepackage[hidelinks]{hyperref}
\setlength{\emergencystretch}{1.2em}
\newtheorem{theorem}{Theorem}
\newcommand{\Z}{\mathbb Z}
\newcommand{\Q}{\mathbb Q}
\newcommand{\F}{\mathbb F}
\newcommand{\Cl}{\operatorname{Cl}}
\newcommand{\Gal}{\operatorname{Gal}}
\newcommand{\len}{\operatorname{length}}
\newcommand{\Sel}{\operatorname{Sel}}
\newcommand{\avref}[2]{\href{run:lean/src/Development/AutomorphicArithmetic/AutomorphicVarieties/#1.lean}{\texttt{#2}}}
\newcommand{\denseref}[2]{\avref{ModularityArgument/ThresholdArguments/DenseThresholdCases/#1}{#2}}
\title{A rank-three congruence bound for the higher \(K\)-groups of \(\Z\)}
\date{}
\author{}
\begin{document}
\maketitle
\vspace{-3.8em}

\begin{abstract}
We extract a fixed-level consequence from the rank-three congruence argument in
Alp\"oge's November 2026 Lean archive. With its arithmetic inputs, we obtain
$\log |K_{4m}(\Z)|\ll(m+1)^7$ and the simultaneous vanishing of a growing
interval of even class-group eigenspaces at every layer of a cyclotomic tower.
The exponent comes from six powers for the coefficient system and one for
the height of a Frobenius trace. This is a mathematical deduction from the
archived inputs. The dependency closure and the deductions here have not
been independently kernel-checked in this study.
\end{abstract}

\section{The conclusions and their input}
For an odd prime $\ell$, write $A_{\ell,n}$ for the $\ell$-primary class group of $\Q(\zeta_{\ell^n})$
and $\omega$ for the Teichm\"uller character of its subgroup of order $\ell-1$
in the cyclotomic Galois group.

\begin{theorem}\label{main}
Using the rank-three comparison and cohomological counting inputs in the
archive, there is an absolute $C$ such that, for all $m\ge1$,
\[
                 \log |K_{4m}(\Z)|\le C(m+1)^7.
\]
For $X\ge3$ and $E\ge2$, the number of primes $\ell\in[X,2X]$ for which
$A_{\ell,n}^{\,\omega^{-e}}\ne0$ for some $n\ge1$ and some even
$2\le e\le E$ is $O((E+1)^8/\log X)$. In particular, outside a set
of primes with convergent reciprocal sum, for every $n\ge1$,
\[
 A_{\ell,n}^{\,\omega^{-e}}=0\qquad
 \text{for every even }2\le e\le(\ell/\log\ell)^{1/8}.
\]
\end{theorem}

The Dirichlet proof in
\href{run:lean/src/Development/AutomorphicArithmetic/CharacterBoxes/MainConclusion.lean}{\texttt{MainConclusion}}
calls the automorphic window. Its proof combines four estimates in
\denseref{WindowBoundOfLiteratureLevel}{the window lemma}. The first, constructed in
\denseref{LinkedUnlE}{LinkedUnlE}, calls both
\denseref{FOneCharU}{FOneCharU} and \denseref{FPropGKU}{FPropGKU}.
The former bounds a Selmer length by valuations of differences of Frobenius
traces on rank-three representations. It passes through a deformation model
and counts inequivalent specializations. The latter puts the necessary
trace differences into one rational integer, simultaneously for all
coefficient primes away from the test prime. The zero-free theorem itself
is unnecessary for what follows.

The convenient specialization is $K=\Q(\sqrt{-3})$, character level $N=2$,
and the trivial character. In
\avref{SmallClassNumbers/ClassNumberThree}{ClassNumberThree},
the class number one of $K$ verifies the geometric hypotheses in
\texttt{gd} for every even $e\ge2$ and every \emph{good} prime $\ell$, meaning
$\ell\ge5$ and $\ell-1\nmid 2e$. At coefficients $\Z_\ell$ the normalization
by the absolute ramification index is one. The archive's
\avref{SmallClassNumbers/CorGKOf}{CorGKOf} already gives an unspecified
polynomial bound for the sum of these lengths weighted by $\log\ell$.
Sections 2 and 3 recover the explicit exponent and identify the full orders.

\newpage
\section{What the Frobenius integer actually measures}
Fix $p=7$ or $11$ and a place $w$ of $K$ over $p$. Put $q=Nw$, so $q=7$ or
$121$, respectively. In the trivial-character specialization the comparison
trace is $q^{-e}+1+q^e$. The construction in
\avref{BuildingsResidueChars/FrobeniusTraceNormNonzeroLogBound}{FrobeniusTraceNormNonzeroLogBound}
forms an integer of the shape
\[
 B_{p,e}=\prod_{s\in D_{p,e}}
 q^e\bigl(t_s-(q^{-e}+1+q^e)\bigr).
 \tag{1}
\]
Here $D_{p,e}$ is a finite indexing set for systems, with possible repetitions
of their traces at $w$. The factors are integral and Galois-permuted. Since
$|t_s|\le3$ and $q\ge7$, none vanishes. Each factor has logarithmic
absolute value at most $2e\log q+1$. The comparison in
\denseref{FPropGKU}{FPropGKU}, as stated in
\avref{UnitaryThreeGroups/KatoUse}{PropGK}, gives for every good $\ell\ne p$
\[
 \len_{\Z_\ell}\Sel_{\ell,e}\ \le\ 2v_\ell(B_{p,e}).
 \tag{2}
\]
Here $\Sel_{\ell,e}$ is defined on page 3. The set $D_{p,e}$ combines the
rank-three CM systems at level $2$ and the rational rank-two systems at
level $2|\operatorname{disc}K|=6$, embedded in rank three.

\paragraph{Keeping weight separate from level.}
In \denseref{TCardK}{TCardK} and \denseref{BrCardP}{BrCardP}, line 30 in
each, the number of systems up to isomorphism is bounded by
$c M^a\dim\Xi$. This follows from the cohomological realization and the
fixed-level cohomology bound. Keep $M$ fixed here. The proof then needs the
dimension of a coefficient system of the indicated infinitesimal character.
For a decreasing integral parameter $b=(b_1,\dots,b_r)$ the single-embedding
factor can be chosen irreducible, of dimension
\[
             \prod_{i<j}\frac{b_i-b_j}{j-i}.
 \tag{3}
\]
Indeed the highest weight in \denseref{TWeylC}{TWeylC} is, with zero-based
indices, $\nu_k=k+1+H-b_{r-1-k}$. Substitution in the Weyl dimension
formula gives (3). A common determinant twist does not affect it. The
irreducible highest weight module supplies the scalar central character used
there. The tensor construction takes one factor per embedding, as recorded in
\avref{GaloisLibrary/WeylCharacterRegularity/ExistsTens}{ExistsTens}.

The actual rank-three parameter is $(e-1,-1,-e-1)$. Formula (3) gives $e^3$
at each of the two embeddings, hence $e^6$. For rank two the parameter
$(e,-e)$ gives $2e$ per embedding. The rational classification in
\denseref{FrobeniusTraceCountBoundOfCmEigensystemCount}{countQ\_alt}
restricts to a fixed finite list of imaginary quadratic fields, at the same
level. Its \avref{GaloisLibrary/CrystallineBaseChange/RootSum}{rootSum}
construction takes a disjoint union, with at most $2^r$ choices of
residue-degree roots and a binary label per system. Thus the rank-two count
is $O(e^2)$. The conversions from eigenvalues to traces preserve the index
cardinalities (lines 348 and 418 onward of the same count file).
\denseref{KeepNumU}{KeepNumU} takes the disjoint union of the two lists.
Consequently $|D_{p,e}|\ll e^6+e^2$ and (1) gives
\[
                   \log|B_{p,e}|\ll(e+1)^7 .
 \tag{4}
\]

\newpage
\section{The Selmer group, including the omitted primes}
Set $D=\Q_\ell/\Z_\ell(-e)$ and $T=\Z_\ell(e+1)$. The definition of
\avref{GaloisLibrary/SelmerTransport/ConjStableShaFam}{shaFam} and its
specialization in \avref{SmallClassNumbers/LenEqZeroIff}{LenEqZeroIff}
identify $\Sel_{\ell,e}\subset H^1(\Q,D)$ as the classes unramified at
every finite prime and locally zero at $2$ and $\ell$. Write
$S=\{\ell,\infty\}$ and $G_S$ for the Galois group of the maximal
extension of $\Q$ unramified outside $S$. For odd $\ell$ and $e>0$,
\[
 \Sel_{\ell,e}=\ker\big(H^1(G_S,D)\longrightarrow H^1(\Q_\ell,D)\big).
 \tag{5}
\]
Indeed $D$ is unramified away from $\ell$, so a cocycle trivial on
inertia there factors through $G_S$. The real place contributes no
positive-degree cohomology. At $2$ an unramified class belongs to
$D/(\mathrm{Frob}_2-1)D=0$, since $2^{-e}-1$ is nonzero and multiplication
by it is surjective on $D$. Local triviality at $\ell$ implies the required
inertia condition there.

Poitou--Tate duality identifies the dual of (5) with the localization kernel
in $H^2(G_S,T)$ \cite[eq.~(22)]{banaszak}. Local Tate duality identifies
the target $H^2(\Q_\ell,T)$ with $H^0(\Q_\ell,D)^\vee$. For good $\ell$
this target is zero: the residual cyclotomic character on $G_{\Q_\ell}$
is surjective and $\ell-1\nmid e$. Quillen--Lichtenbaum then gives
\cite[Thm.~70]{weibel}
\[
       \len_{\Z_\ell}\Sel_{\ell,e}
       =\len_{\Z_\ell}H^2(G_S,T)
       =v_\ell\bigl(|K_{2e}(\Z)|\bigr).
 \tag{6}
\]
Thus (2) bounds prime-power multiplicities.

It remains to account for primes omitted by (2). For every odd $\ell$ the
standard mod-$\ell$ comparison \cite[Thm.~5.1]{jkpsz} gives
\[
 K_{2e}(\Z)/\ell \ \simeq\
 \bigl(\Cl(\Q(\zeta_\ell))/\ell\bigr)^{\omega^{-e}}. \tag{7}
\]
When $\ell-1\mid2e$ the character on the right is trivial or quadratic.
The trivial part comes from $\Cl(\Q)=0$. In the quadratic case
$e$ is even, hence $\ell\equiv1\bmod4$ and the eigenspace comes from
$\Q(\sqrt{\ell})$. In both cases the descent assertion follows from the
norm and extension of ideals, since the extension degree is prime to $\ell$.
Minkowski supplies an integral ideal of norm at most $\sqrt{\ell}/2$ in
each real quadratic ideal class. There are at most $\tau(j)\le j$ ideals
of norm $j$, so $h(\Q(\sqrt\ell))<\ell$. The quadratic part is zero as
well. This includes all the needed cases at $3$ through the trivial
character. In degree $2e$ divisible by four the $2$-part is zero
\cite{rognes}. Finally $7$ and $11$ are regular primes, as is seen from
$B_2,B_4,B_6,B_8$ and Kummer's criterion, so (7) rules out both.
The groups in question are finite. Combining (2) and (6) at every
possible divisor gives the more precise statement
\[
        |K_{2e}(\Z)|\ \mid\ \gcd(|B_{7,e}|,|B_{11,e}|)^2
        \qquad (e\ge2\text{ even}). \tag{8}
\]
Taking $e=2m$ in (4) and (8) proves the first assertion. The construction
gives existence and a height bound for the integers. Numerical values of
the trace products have not been computed here.

\newpage
\section{A growing interval throughout the cyclotomic tower}
By (7), $A_{\ell,1}^{\omega^{-e}}\ne0$ exactly when $\ell$ divides
$|K_{2e}(\Z)|$. The unique prime over $\ell$ is totally ramified throughout
the cyclotomic $\Z_\ell$-extension. The norm control theorem therefore
identifies the $\Gamma_n$-coinvariants of the unramified Iwasawa module
$X_\infty$ with $A_{\ell,n}$, where $\Gamma_n$ fixes $\Q(\zeta_{\ell^n})$
\cite[chapter~3]{sharifi}. All maps respect the prime-to-$\ell$ character
action. In particular if the $\omega^{-e}$ part at the bottom is zero,
the same part of $X_\infty$ has zero $\Gamma_1$-coinvariants.
Nakayama's lemma over $\Z_\ell[[\Gamma_1]]$ gives its vanishing, hence
vanishing at every level. The usual finite generation of the unramified
Iwasawa module suffices for this application.

Every exceptional prime in $[X,2X]$ thus divides the single integer
\[
                    \prod_{\substack{2\le e\le E\\ e\ {\rm even}}}
                           |K_{2e}(\Z)|.
\]
The logarithm of this integer is $O((E+1)^8)$ by (4) and (8). Each
distinct exceptional prime costs at least $\log X$, proving the claimed
count. For the displayed window of Theorem~\ref{main}, take
$E=(2X/\log X)^{1/8}$ on each dyadic interval. The exceptional count is
$O(X/(\log X)^2)$, giving a convergent sum of reciprocals and relative
natural density zero among primes. For an individual prime, (4) and (8)
also force vanishing whenever $\log\ell>C'(e+1)^7$.

\paragraph{Which characters this reaches.}
The exponent $-e$ means the even characters with representatives just below
$\ell-1$. Chen et al.\ obtain a density-one result for the even characters
with small \emph{positive} exponents, up to $\sqrt{\ell}/(\log\ell)^\alpha$
for $\alpha>1/2$ \cite{chen}. The two intervals can be imposed on the
same density-one set of primes. Soul\'e treated the negative-exponent
direction for individual primes, with a threshold
$\log\ell>n^{224n^4}$ for $n=e+1$ in his notation \cite{soule}.
The source-level exponent $7$ above is an upper estimate. Its optimality
and the priority of these deductions are open in this note.

\small
\begin{thebibliography}{10}
\bibitem{alp} L. Alp\"oge, \href{https://alpo.ge/remark.pdf}{announcement and Lean archive}, 20 November 2026. The typewriter links above point into the supplied local source tree.
\bibitem{banaszak} G. Banaszak, \href{https://arxiv.org/abs/1208.2137}{Wild kernels and divisibility in \(K\)-groups of global fields}.
\bibitem{weibel} C. Weibel, \href{https://sites.math.rutgers.edu/~weibel/papers-dir/KZsurvey-published.pdf}{Algebraic \(K\)-theory of rings of integers in local and global fields}.
\bibitem{jkpsz} B. Jordan, Z. Klagsbrun, B. Poonen, C. Skinner and Y. Zaytman, \href{https://arxiv.org/abs/1703.00108}{Statistics of \(K\)-groups modulo \(p\) for the ring of integers of a varying quadratic number field}.
\bibitem{rognes} J. Rognes and C. Weibel, \href{https://sites.math.rutgers.edu/~weibel/archive/papers-dir/RognesWeibel.pdf}{Two-primary algebraic \(K\)-theory of rings of integers in number fields}.
\bibitem{sharifi} R. Sharifi, \href{https://math.ucla.edu/~sharifi/notes/iwasawa-ch03.html}{Iwasawa Theory}, chapter 3, limits of class groups.
\bibitem{chen} E. Chen et al., \href{https://arxiv.org/abs/2602.05090}{Almost all primes are partially regular}.
\bibitem{soule} C. Soul\'e, \href{https://arxiv.org/abs/math/9812171}{Perfect forms and the Vandiver conjecture}.
\end{thebibliography}
\end{document}
